%==========================================================================================
\section{\texttt{xrf\_region}: surface region of a map}\label{sec:region}
%==========================================================================================

\subsection{Overview}

A thin film covers the whole map, and every pixel belongs to the same flat surface. A single
crystal glued on the holder does not: the map shows a central flat facet, neighbouring facets
at other tilts, the glue and the holder. A facet tilted by $\delta$ changes the incidence and
take-off angles by up to $\delta$, and with them the absorption paths of every line
(section~\ref{sec:matrix}); its intensities and line ratios therefore differ from those of the
central facet even at the same composition. Only the facet whose orientation is known (on
BM32, the plane fitted to align the \fn{sx}/\fn{sy} frame) belongs in the average spectrum and
the quantification.

The module describes that facet as a \fn{SurfaceRegion}: polygons that are included, minus
polygons that are excluded (an inclusion, a dust grain, a crack). The region gives a boolean
mask over the pixels, which the notebook applies to the average spectrum, the pile-up template,
the per-segment fits (\fn{fit_segments}) and the map solve (\fn{fit_map}, argument
\fn{mask}). The vertices can be typed in, or drawn on the map with \fn{RegionSelector}, and
saved to a small JSON file.

\begin{hyp}
\begin{itemize}[nosep,leftmargin=1em]
\item A pixel belongs to the region when its \emph{centre} (the recorded motor positions of
  the point) is inside; the beam footprint is not considered, so pixels along an edge mix two
  facets whichever side they fall on. Drawing the polygon slightly inside the facet avoids them.
\item Motor units are the recorded positions of the fast and slow motors of the mesh, read
  point by point: they do not depend on the grid and remain valid for another map of the same
  facet in the same frame. They are only meaningful in the frame they were drawn in, so the
  motor names are stored and checked.
\item Points exactly on an edge may fall on either side (\fn{Path.contains_points}). In pixel
  units with integer vertices, the pixel centres lie on the edges: shift the vertices by half
  a pixel to make the choice explicit.
\end{itemize}
\end{hyp}

\pycode{xrf_region}{_header}

\subsection{Polygons and coordinates}

\api[\texttt{polygon\_area}, \texttt{points\_in\_polygon}]{\texttt{polygon\_area(vertices)}, \texttt{points\_in\_polygon(x, y, vertices)}}
\fn{_as_polygon} checks and normalises a list of vertices to an $(n,2)$ array (a repeated
closing vertex is dropped, at least three distinct vertices are required).
\fn{polygon_area} is the shoelace formula,
$A = \tfrac12\left|\sum_i (x_i y_{i+1} - y_i x_{i+1})\right|$.
\fn{points_in_polygon} returns, for each point, whether it is inside the polygon, from
\fn{matplotlib.path.Path.contains_points} (even--odd rule).
\pycode{xrf_region}{_as_polygon}
\pycode{xrf_region}{polygon_area}
\pycode{xrf_region}{points_in_polygon}

\api[\texttt{scan\_coordinates}]{\texttt{scan\_coordinates(scan, units="motor")}}
The $(x, y)$ = (fast, slow) coordinates of every pixel of a mesh scan, as flat arrays in
acquisition order. \texttt{"motor"}: the recorded positions of the two motors (they must have
been read, which \fn{BlissMeshScan.read_data} does by default); \texttt{"pixel"}: the
(column, row) indices on the \fn{scan_shape} grid.
\pycode{xrf_region}{scan_coordinates}

\api[\texttt{facet\_plane}]{\texttt{facet\_plane(scan, motors=("xech", "yech", "zech"), mask=None)}}
The plane $z = ax + by + c$ followed by the sample stage during a map, by least squares on the
recorded positions of the three motors. On BM32, the map of a crystal is scanned in the
\fn{sx}/\fn{sy} frame, which is built (\fn{initmat}) from a plane fitted on the facet before
the map; the stage then follows that plane point by point and the fit returns the facet
orientation in the stage frame (residual of a few nm), for \fn{LabGeometry.facet}. Returns
\fn{a}, \fn{b}, \fn{c}, the residual \fn{rms_um}, the tilt \fn{tilt_deg} $= \arctan\sqrt{a^2+b^2}$
and the number of points.
\pycode{xrf_region}{facet_plane}

\subsection{The region}

\api[\texttt{SurfaceRegion}]{\texttt{SurfaceRegion(include=[], exclude=[], units="motor", motors=None, name="", note="")}}
\begin{args}
\item[include] polygons whose pixels are kept; empty means the whole map.
\item[exclude] polygons whose pixels are removed from the included ones.
\item[units] \texttt{"motor"} or \texttt{"pixel"} (see above).
\item[motors] the (fast, slow) motor names, checked against the scan by \fn{scan_mask}.
\item[name, note] free text, kept in the JSON file.
\end{args}
\pycode{xrf_region}{SurfaceRegion}
\pycode{xrf_region}{SurfaceRegion.__post_init__}

\api[constructors]{\texttt{SurfaceRegion.from\_vertices(vertices, units="motor", exclude=(), **kw)}, \texttt{SurfaceRegion.coerce(region)}}
\fn{from_vertices} makes a region of one polygon, the usual case of one facet. \fn{coerce}
accepts what a notebook configuration naturally contains: \texttt{None} (the whole map), a
\fn{SurfaceRegion}, the path of a JSON file, a dictionary of keyword arguments, or a list of
vertices in motor units.
\pycode{xrf_region}{SurfaceRegion.from_vertices}
\pycode{xrf_region}{SurfaceRegion.coerce}

\api[\texttt{mask}, \texttt{scan\_mask}]{\texttt{is\_full}, \texttt{mask(x, y)}, \texttt{scan\_mask(scan)}, \texttt{area()}}
\fn{mask} is true for the points inside any included polygon and inside no excluded one;
\fn{scan_mask} applies it to the pixels of a scan (a flat $(n_\mathrm{pixels},)$ array, all
true for the whole map) and warns when the motors differ from those the region was drawn with,
or when no pixel is selected. \fn{area} is the included minus the excluded polygon areas
(overlaps are not resolved), in mm$^2$ or pixels.
\pycode{xrf_region}{SurfaceRegion.is_full}
\pycode{xrf_region}{SurfaceRegion.mask}
\pycode{xrf_region}{SurfaceRegion.scan_mask}
\pycode{xrf_region}{SurfaceRegion.area}
\pycode{xrf_region}{SurfaceRegion.__str__}

\api[\texttt{plot}]{\texttt{plot(ax, origin=(0, 0), scale=1.0, include\_kw=None, exclude\_kw=None)}}
Draws the outlines (included: solid white, excluded: dashed red) at
$(\mathbf v - \mathbf o)\,s$, so that they can be laid over a map drawn in other units, e.g.
\textmu m from the corner of the map with \fn{origin} the minimum motor positions and
$s = 1000$.
\pycode{xrf_region}{SurfaceRegion.plot}

\api[input/output]{\texttt{to\_dict()}, \texttt{save(path, **metadata)}, \texttt{SurfaceRegion.load(path)}, \texttt{to\_code(precision=5)}}
The region is saved as JSON (units, motors, vertices; extra metadata such as the scan file are
stored alongside and ignored on reading). \fn{to_code} prints a Python expression recreating
it, to paste into a notebook configuration.
\pycode{xrf_region}{SurfaceRegion.to_dict}
\pycode{xrf_region}{SurfaceRegion.save}
\pycode{xrf_region}{SurfaceRegion.load}
\pycode{xrf_region}{SurfaceRegion.to_code}

\subsection{Drawing a region}

\api[\texttt{map\_image}]{\texttt{map\_image(corrected, calibration=None, window\_ev=None)}}
The per-pixel image to draw on: the sum of the normalised spectra, over an energy window if
given. A strong line of the crystal shows the facets best, since their intensity depends on
the local angles.
\pycode{xrf_region}{map_image}

\api[\texttt{RegionSelector}]{\texttt{RegionSelector(scan, image, units="motor", region=None, name="", log=False, cmap="viridis", figsize=(7.0, 6.2))}}
An interactive map built on Matplotlib's \fn{PolygonSelector}. It needs an interactive backend:
\texttt{\%matplotlib widget} (ipympl) in Jupyter, where buttons are added with ipywidgets, or a
Qt/Tk window in a script. The user clicks the vertices of the facet and closes the polygon on
the first vertex; vertices can then be dragged, or the whole polygon with \texttt{shift}.
\textbf{Add facet} (key \texttt{a}) adds it to the region, \textbf{Add exclusion} (key
\texttt{x}) to the exclusions, \textbf{Undo} (\texttt{u}) removes the last one, \textbf{Clear
all} returns to the whole map. Pixels outside the current region are veiled in grey and the
number of selected pixels is shown, so what will be averaged is visible while drawing.
\fn{region} and \fn{mask} return the result so far. With a non-interactive backend (inline)
the map is still shown and a warning explains how to switch.
\pycode{xrf_region}{RegionSelector}
\pycode{xrf_region}{RegionSelector.__init__}
\pycode{xrf_region}{RegionSelector._show}
\pycode{xrf_region}{RegionSelector._title}
\pycode{xrf_region}{RegionSelector._update_status}
\pycode{xrf_region}{RegionSelector._redraw}
\pycode{xrf_region}{RegionSelector._on_polygon}
\pycode{xrf_region}{RegionSelector._on_key}
\pycode{xrf_region}{RegionSelector.add}
\pycode{xrf_region}{RegionSelector.undo}
\pycode{xrf_region}{RegionSelector.clear}
\pycode{xrf_region}{RegionSelector.region}
\pycode{xrf_region}{RegionSelector.mask}
